% =====================================================================
% chapters/06_Non_tabulated.tex — LaTeX rendering of ../sections/06_non_tabulated.en.md (2026-09-23 12:55:30)
% Section 6 of the SHORT paper. \input from main.tex, after 05_Results.tex.
% This file DEFINES \label{sec:nontabulated}, referenced by the reading map of Section 1.4,
%   and \label{sec:mechanism}, referenced from Section 7.1.
% Figure: one, \columnwidth (../PLAN.md § 9, figure 5, one column), caption BELOW and a
%   mandatory \Description{}.
%   FILE: images/ch7_propension_quotidienne.png — REGENERATION PENDING: its title, its axis
%   labels, its legend ("témoin", "exposé") and its annotations ("avarie", "le récit quitte le
%   contexte") are still in French. Every figure of the paper is composed in English.
% Table: one, tab:offered, \columnwidth, caption ABOVE, booktabs. ../PLAN.md § 9 caps the
%   paper at three tables; the review (item N2) numbers this one as Table 4 for consistency.
%   IF THE COMPILED PDF RUNS PAST EIGHT PAGES, this table is the first thing to drop: the
%   Markdown holds the replacement sentence, ready to paste —
%   "The car is taken on 94\,\%, 32\,\% then 88\,\% of the trips where it is offered, against
%   90, 92 and 87\,\% for the control agent."
% No citation in this section.
% Derived file: regenerate after every validated pass of ../sections/06_non_tabulated.en.md.
% =====================================================================

\section{The non-tabulated regime: a single-agent event traced to the decision}
\label{sec:nontabulated}

This section follows one event that no variable carries, in one agent, from its entry into
memory to the decisions it changed.

\subsection{What the 21 variables do not carry}
\label{sec:variables-gap}

None of the 21 variables records what a person went through yesterday, or what they read
this morning. A decision-maker that receives those inputs alone can answer nothing else. Its
prediction on the morning after an incident equals the one before, by identity rather than
by measurement.

\subsection{The mechanism}
\label{sec:mechanism}

An event enters the agent's memory at one of two moments. Something lived enters after the
trip that produced it, and something read enters at waking. At entry the agent rates the
event on a five-step severity scale, and gives it a valence, suffered or welcome. The rating
is the agent's own, read against its profile, so one incident weighs differently on two
agents. The severity alone sets how long the memory stays in service.

A rated event travels beyond its owner, through the household in the evening.

An effect ends in one of two ways, which measurement tells apart. It ends by
contradiction when the agent makes the trip again and nothing recurs, and by wear when the
memory stops being served.

Two links of this chain are text generation. Writing the trace of what has just
happened is one, and maintaining the belief drawn from it is the other. The language model
rates the severity as well, and that link is not text generation. The rating is one level in
a closed grid, which a typed classifier can return. Whether it could also produce the two
written links is open, and one bench arm would settle it.

\subsection{The setup}
\label{sec:setup}

We play one agent twice, exposed and control, and read the gap between the two runs. The
incident is an engine failure that imposes a half-hour delay on a car trip. Both runs keep
memory enabled. The car remains offered the next morning, so that what we measure here is an
arbitration. \textbf{[TBC --- A campaign now under way]}

Opinions are collected outside any decision, on six declared criteria, at four milestones.
The criterion on the environment, which no incident targets, serves as the control question.

\subsection{Results}
\label{sec:traced-results}

Although the car is offered as often as before the failure, it stops being taken. The
exposed agent takes it nine times in ten before the failure, and three times in ten after.
The control agent holds its rate throughout.

\begin{table}[tbp]
  \caption{Car taken when the car is offered, exposed and control agent.}
  \label{tab:offered}
  \begin{tabular}{@{}lrr@{}}
    \toprule
    Period & Exposed agent & Control agent \\
    \midrule
    Before the failure & 94\,\% (34/36) & 90\,\% (28/31) \\
    After the failure  & 32\,\% (12/38) & 92\,\% (45/49) \\
    After the return   & 88\,\% (35/40) & 87\,\% (40/46) \\
    \bottomrule
  \end{tabular}
\end{table}

The daily propensity to the car falls on the day of the failure and returns
\textbf{fifteen days} later (Figure~\ref{fig:propensity}). It drops from 90\,\% on the eve to 40\,\% that day,
then oscillates between 5 and 52\,\%. It rejoins the band of the control agent, which never
leaves 65 to 90\,\%, on 15 April.

\begin{figure}[tbp]
  \centering
  \includegraphics[width=\columnwidth]{images/ch7_propension_quotidienne.png}
  \caption{The effect lasts exactly as long as the memory stays in the prompt.}
  \Description{Two daily time series of the probability the agent assigns to the car, in per cent, over about six weeks of simulated days from mid-March to late April. The exposed agent and the control run together above 65 per cent until the day of the engine failure, marked by a dashed vertical line. The exposed series then collapses to between 5 and 52 per cent for fifteen days, while the control keeps oscillating between 65 and 90 per cent. A dotted vertical line marks the day the account of the failure leaves the decision context, and the exposed series rejoins the control band two days later, on 15 April, and stays in it.}
  \label{fig:propensity}
\end{figure}

The four routes by which the past reaches a decision can be told apart, and the trace names
the one that carried this effect. The account of the failure appears in the block of what has
changed recently, in 75 of 376 decision prompts. It stays there from the day of the
failure to the \textbf{fifteenth} day after. That span follows from the severity.
\textbf{[TBC --- The campaign under way reads what the three other routes carry.]}

Measurement separates the two ways an effect can end, and this one ended by wear. The agent
took the car twelve times during the period. \textbf{[TBC --- at this time, no contradiction
of the belief born of the failure. Add an experiment?]} The account leaves the decision prompts after 13 April, and
the propensity rejoins the control band on 15 April. Two independent sources give those
dates, the text sent to the model and the probabilities it returned.

Declared opinions drop after the failure and come back. Five of the six criteria on the car
fall at the milestone that follows it, safety from 8 to 3 on a scale of ten. The five have
recovered their starting value at the two later milestones. The environment criterion stays
at 3 throughout, and the control agent moves on none of the six. Both agents
drift on the four other modes without a common pattern, leaving the car as the only mode
that separates them.

\textbf{A second experiment carries the same channel to information that is only read. Now
under way, it serves five dated Toulouse press articles against twenty signs written before
any call to the model (Appendix E).}

The path traced here rests on one agent and one event. We claim the path rather than its
amplitude.
